In 1983 bracht Richard Stallman\index{Stallman, Richard} een nieuw project in de wereld genaamd het \index{GNU}GNU Project, waarbij GNU staat
voor GNU's not Unix! Dit om aan te geven dat het gebaseerd is op Unix maar geen Unix componenten bevat en er dus geen licentie van AT\&T (Bell Labs) nodid was. Het systeem is \textquote{free software} als in dat je niet afhankelijk bent van anderen om het systeem te bestuderen en/of te wijzigen. GNU is een \textquote{from scratch} geschreven systeem dat volledig open source is. From-scratch betekent dat alle code opnieuw geschreven is en Open Source wil zeggen dat van het hele systeem alle broncode openbaar
beschikbaar is, hierdoor kan het gebruikt worden als studiemateriaal. Het heeft tevens het voordeel dat iedereen fouten
uit de software kan halen en aan de software kan bijdragen om het beter te maken.

Om de vrijheid van software ontwikkeling en bestudering mogelijk te maken moest er een licentie komen waarin deze rechten waren opgenomen. Deze licentie werd de \index{General Public License}General Public License, afgekort de \index{GPL}GPL. De GPL zegt dat je de sofware mag kopie\"eren, bestuderen en aanpassen. Het enige dat je niet mag is de licentie aanpassen.

De software van het GNU project werd door veel verschillende distributies, commercieel en niet commercieel, gebruikt om software uit het oorspronkelijke Unix geheel of gedeeltelijk
te vervangen. Missend onderdeel in het GNU-project was jarenlang een kernel. Het stukje software dat ligt tussen de gebruikersinterface en de hardware.

